\section{E. Small deformations and the comb construction}

\begin{claim}{E1.1}{Proved}{Stability of the pencil for the standard family}
Every sufficiently small deformation of the standard $X_c$ carries a
holomorphic map to $\Pone$ with general fibre a complex two-torus. Its
complete anticanonical system $|-2K|$ is basepoint-free of dimension one
and recovers that map, and its algebraic dimension is one.
\end{claim}
\begin{proof}
D3.1 identifies the semiuniversal germ with the actual $c$-parameter
germ. Thus every sufficiently small deformation is isomorphic to a
member $X_{c'}$ with $\operatorname{Im}c'<-M$. The bundle identity
$K_X=\bar f^*\OO_{\mathcal B}(-6)=\OO_X(-2S_2)$ in D1.2 holds
relatively in this family: all root divisors and every determinant
transition used in its proof vary holomorphically, and their characters
are constant. Consequently $-2K=\bar f^*\OO(12)=f^*\OO_B(1)$
on every member, with a relative line-bundle isomorphism. Connected
fibres imply $f_*\OO_X=\OO_B$, so the two base sections pull back to
the entire two-dimensional section space. They have no common zero.
The resulting map is the constructed $f$ followed by a projective
coordinate change. Its general fibre is a torus by the same local
quotient construction.

We spell out the algebraic-dimension reduction used here. On a
manifold with $H^2(X,\R)=0$, a divisor restricts to a topologically
trivial line bundle on any regular torus fibre. A nonzero effective
divisor on a compact torus has positive degree against a Kahler form,
whereas a topologically trivial class has degree zero. Hence an
effective divisor on $X$ cannot meet a regular fibre in a proper
nonempty divisor: its restriction is either the whole fibre or empty.
The zero and pole divisors of a meromorphic function are therefore
vertical. On a regular fibre where neither contains the fibre, the
restricted function has no pole and is constant. The function hence
descends meromorphically to the base on a dense open set, and the
proper connected-fibre map extends the descent across the finitely
many special values. Thus the meromorphic field is $\C(t)$ and
$a(X)=1$. This argument applies to each small member and does not
appeal to an invalid vanishing theorem.

The conclusion is stronger than what semicontinuity alone gives.
For the original member, projection formula gives
$h^\bullet(f^*\OO(1))=(2,3,1,0)$; the nonzero degree-one group means
that a generic line-bundle extension argument would not establish
stability of its two sections. Here stability follows from the
computed semiuniversal family and its relative root divisors.
\end{proof}

\begin{claim}{E2.1}{Known}{The structure sheaves of the comb members}
Let $X_m$ be any sufficiently small smooth period section obtained by
the $m$-tooth construction of
\citet[Theorem 1.1]{ViaclovskyCombs2026}, with
$m\ge12$, fixed limiting data at the three distinguished values, and
the standard logarithmic-transformation data. Write
$f_m:X_m\to B$. Then
\[
 R^1f_{m*}\OO_{X_m}=\OO(a_m)\oplus\OO(b_m),\quad
 a_m,b_m\ge0,\quad a_m+b_m=2m-1,
\]
\[
 R^2f_{m*}\OO_{X_m}=\OO(2m-1),\qquad
 (h^{0,0},h^{0,1},h^{0,2},h^{0,3})=(1,2m+1,2m,0).      \tag{E.1}
\]
The individual summands are not specified by $m$ in the construction.
\end{claim}
\begin{proof}
The paper's equation (7.5) proves that attaching and smoothing one
twistor conic increases the determinant degree of its quotient bundle
by two. Its Proposition 9.1 identifies the initial degree as $-1$,
so the smoothed quotient has degree $2m-1$. The geometric gluing
comparison in its Section 6.3 identifies that quotient with
$R^1f_{m*}\OO_{X_m}$; the fixed limiting periods preserve the same
cyclic representations at the finite points and the same
normalization complex at the cusp. Thus its Theorem 2.3 applies to
these actual full scheme-theoretic fibres and gives the cup-product
isomorphism $\det R^1f_{m*}\OO=R^2f_{m*}\OO$.
Its Proposition 3.2, whose hypotheses are connected fibres, torus
general fibre and the topology of $S^6$, gives the two possible
splitting alternatives: both degrees nonnegative, or $(-1,0)$.
The positive determinant degree excludes the latter. We have now
derived the stated sheaves from the source comparisons in precisely
the form needed here. Leray on $B$ gives $h^{0,1}=a_m+b_m+2$,
$h^{0,2}=a_m+b_m+1$ and $h^{0,3}=0$, proving (E.1).
The theorem is used for the smoothed section, not for the singular
comb, which is not a small deformation of $X_0$.
\end{proof}

\begin{claim}{E2.2}{Proved}{Canonical bundle, exact tangent cohomology and Hodge growth on the combs}
For the same $X_m$, the root lift $\bar f_m:X_m\to\mathcal B$ has
vertical bundle $E_m$ and $G_m=R^1\bar f_{m*}\OO_{X_m}$ satisfying
\[
 0\longrightarrow\OO\longrightarrow E_m
       \longrightarrow\OO(24m-1)\longrightarrow0,
 \qquad
 0\longrightarrow\OO\longrightarrow G_m
       \longrightarrow\OO(24m-1)\longrightarrow0.      \tag{E.2}
\]
The canonical bundle and all positive-degree holomorphic forms obey
\[
 K_{X_m}=\bar f_m^*\OO(-24m-6),\qquad
 h^{1,0}=h^{2,0}=h^{3,0}=0.                            \tag{E.3}
\]
Moreover
\[
 \boxed{h^\bullet(T_{X_m})=(2m+1,8m+1,6m+1,0).}       \tag{E.4}
\]
More precisely, writing $\mathcal L=\OO_{\mathcal B}(-5)$,
\[
 \begin{split}
 f_{m*}T_{X_m}&=\rho_*E_m,\\
 R^1f_{m*}T_{X_m}&=
  \operatorname{coker}\{\rho_*\mathcal L\xrightarrow{\delta_0}
                                      \rho_*(E_m\otimes G_m)\},\\
 R^2f_{m*}T_{X_m}&\simeq\OO_B(2m)\oplus Q_m,
 \qquad\operatorname{length}Q_m=4m,
 \end{split}                                        \tag{E.4a}
\]
where $Q_m=\operatorname{coker}\{\rho_*(\mathcal L\otimes G_m)
\xrightarrow{\delta_1}\rho_*(E_m\otimes\det G_m)\}$ has no support
at the cusp. The splitting in the last line need not be canonical.
There is a rank-three vector bundle $K_m$ on $B$ for which
\[
 \boxed{h^{1,1}(X_m)=4m+1+h^0(K_m),\qquad
 h^{1,2}(X_m)=4m+h^0(K_m).}                            \tag{E.5}
\]
In particular both remaining middle Hodge numbers grow at least
linearly with $m$.
\end{claim}
\begin{proof}
The construction retains the integral local system and the limiting
period planes. The invariant real homology class $\widehat\delta$
therefore gives a vertical translation field on every regular torus.
It is nowhere zero there, since the real period map is an isomorphism.
At a finite root point the same reasoning applies on the smooth
cover, and at the cusp its nonzero angular class gives a nonzero
logarithmic vertical vector. The local vertical transitions, which
depend only on the base exactly as in D1.2, define $E_m$ and give its
trivial subbundle. The invariant real cohomology class $\gamma$ gives
a nowhere-zero section of $G_m$, by the same smooth-fibre and cusp
comparison in D1.2.

The determinant of $G_m$ has the same two stabilizer characters as
that of $G_0$: invariant exponents $2$ at $p_1$ and $1$ at $p_2$.
Its coarse determinant has degree $2m-1$ by E2.1. The explicit descent
formula in D1.1 therefore identifies
$\det G_m=\OO(12(2m-1)+8+3)=\OO(24m-1)$.
The rank-four flat bundle on the root base has trivial determinant:
the integral monodromies have determinant one, the finite-root charts
have no residual ramification, and the logarithmic cusp monodromy
has nilpotent residue of trace zero. Its flat volume section thus
extends without a zero or pole. The exact Hodge-plane sequence
$0\to E_m^*\to\mathcal H\to G_m\to0$ gives
$\det E_m=\det G_m$. Taking quotients by the two nonvanishing
sections proves (E.2). The determinant of the log tangent sequence
$0\to\bar f_m^*E_m\to T_{X_m}(-\log W)
\to\bar f_m^*\OO(-5)\to0$ then gives (E.3)'s canonical formula.

The logarithmic Kodaira--Spencer map is invertible near the cusp.
This does not follow merely from fixing a limiting plane:
Viaclovsky's Proposition 6.1 gives the actual local period expression
$(\log q)I/(2\pi i)+H_m(q)$ with $H_m$ holomorphic. Its logarithmic
derivative has invertible leading matrix $I/(2\pi i)$, or
$B_0/(2\pi i)$ after the retained marking change. Thus the map is
generically invertible. The direct image of the logarithmic
cotangent sequence consequently has kernel exactly $\OO(5)$.
It has no section by D1.1, so $h^{1,0}=0$. To treat two-forms, take
the second exterior power of that sequence. Its possible global
sections lie in
$H^0(E_m^*\otimes\OO(5))$ or
$H^0((\det E_m)^{-1})$.
The latter is zero, having weight $1-24m<0$. A morphism
$E_m\to\OO(5)$ vanishes on the subline $\OO$, since
$H^0(\OO(5))=0$, and therefore factors through $\OO(24m-1)$.
The resulting section would have weight $6-24m<0$, so it too
vanishes. Absolute two-forms inject into logarithmic ones, which
proves $h^{2,0}=0$. The negative canonical weight proves $h^{3,0}=0$.

The first connecting map from the horizontal line to
$E_m\otimes G_m$ is nonzero near the cusp, so it is injective as a
sheaf map. The same normal quotient (D.10) has no degree-zero direct
image, giving $f_{m*}T_{X_m}=\rho_*E_m$.
Exact invariants applied to (E.2) give, for each of $V=E_m,G_m$,
\[
 0\longrightarrow\OO_B\longrightarrow\rho_*V
 \longrightarrow\OO_B(2m-1)\longrightarrow0.
\]
These bundles have vanishing $H^1$ and $h^0=2m+1$.
They also split into two nonnegative-degree lines: if a splitting
had one negative summand, the section of its saturated $\OO_B$
subline would have to lie in the positive summand. Every nonzero
section of that summand has a zero, contradicting saturation.

The natural map
\[
 (\rho_*E_m)\otimes(\rho_*G_m)
              \longrightarrow\rho_*(E_m\otimes G_m)
\]
is injective, with length-one cokernel at $p_1$ and no other
cokernel. Indeed both bundles have invariant exponents $(0,2)$ at
$p_1$. The product of their exponent-two generators is $s^4$,
whereas the product representation has invariant generator $s$;
the comparison on that one summand is multiplication by $t=s^3$.
All other exponents add without a carry. At $p_2$ the exponents
are $(0,1)$, and no sum reaches four. At the cusp there is no
stabilizer. The source is a sum of nonnegative-degree line bundles,
so its $H^1$ is zero. The length-one upper modification consequently
has $H^1=0$, rank four and degree
$2(2m-1)+2(2m-1)+1=8m-3$; it has $8m+1$ global sections.

In the direct-image sequence of the log tangent extension, the
second cup map
$\delta_1:\mathcal L\otimes G_m\to E_m\otimes\det G_m$
is generically invertible, by the invertible cusp Kodaira--Spencer
matrix. It is therefore injective as a sheaf map. The first cup map
$\delta_0$ is injective for the same reason. Consequently
$R^1\bar f_{m*}T(-\log W)=\operatorname{coker}\delta_0$.
The normal quotient $\nu_*K_P$ in (D.10) has zero direct images in
degrees zero and one, so the same formula holds for $R^1T$.
Since $\rho_*\mathcal L=\OO_B(-1)$ is acyclic, its cokernel has
$h^0=8m+1$ and $h^1=0$. Leray and $H^1(\rho_*E_m)=0$ now give
the exact value $h^1(T)=8m+1$.

For completeness the degree-two calculation gives (E.4) directly,
with the full torsion correction retained. The coarse degrees of
$\mathcal L\otimes G_m$ and $E_m\otimes\det G_m$ are $2m-2$
and $6m-2$. Their invariant exponents at $p_1$ are respectively
$(1,0)$ and $(2,1)$, and at $p_2$ both are $(1,2)$.
Their root determinant weights are $24m-11$ and $72m-3$;
subtracting the character contributions and dividing by twelve
gives exactly these coarse degrees. Thus $Q_m$ has length $4m$.
It has no cusp contribution, where the cup map is invertible.
On the root base there is an exact sequence
\[
 0\longrightarrow\operatorname{coker}\delta_1
 \longrightarrow R^2\bar f_{m*}T(-\log W)
 \longrightarrow\mathcal L\otimes\det G_m
                 =\OO(24m-6)\longrightarrow0.
\]
It splits, since the extension group is the first cohomology of a
zero-dimensional sheaf. The cusp normal quotient then raises the
degree of the free line by $p_0$, exactly as in D1.3. The local
fibre, its normalization and the logarithmic residue are unchanged;
in particular $h^2(T|_W)=1$ still excludes a split torsion summand at
the cusp. Hence the new free line is $\OO(24m+6)$, which descends to
$\OO_B(2m)$, proving (E.4a). Leray gives
$h^2(T)=(2m+1)+4m=6m+1$ and $h^3(T)=0$.
As an independent check, these values satisfy the full tangent
index $\chi(T)=1$. The source family of Viaclovsky's Proposition 8.3
also gives
\[
 8m-8\le\dim\operatorname{Def}(X_m)\le8m+1;
\]
the lower bound concerns the actual classifying-family dimension,
while the upper bound is the tangent dimension now computed.

We give the bundle specifying the remaining Hodge number. Set
\[
 K_m=\ker\left\{\rho_*(E_m^*\otimes G_m)
  \xrightarrow{\ \kappa_1\ }
       \rho_*(\OO(5)\otimes\det G_m)\right\}.         \tag{E.6}
\]
The map is cup product with the actual logarithmic
Kodaira--Spencer class, in the category of coherent sheaves on $B$.
It is generically onto, so the kernel is a rank-three vector bundle.
The preceding connecting map
\[
 \rho_*E_m^*\xrightarrow{\ \kappa_0\ }
                 \rho_*(\OO(5)\otimes G_m)
\]
is injective. The source and target have degrees $-2m-1$ and
$2m-1$. To check these degrees, use their fixed local character
exponents $(0,1)$ versus $(2,1)$ at the order-three point and
$(0,3)$ versus $(3,0)$ at the order-four point, together with the
root determinant degrees. The invariant determinant comparison in
D1.1 gives precisely the displayed integers. Its cokernel $A_m$
therefore has length $4m$, irrespective of how its zeros are
distributed among regular and special fibres.

Writing $\mathcal A_m=\Omega^1_{X_m}(\log W)$, the direct-image
sequence gives
$0\to A_m\to R^1f_{m*}\mathcal A_m\to K_m\to0$.
The residue sequence
$0\to\Omega^1_{X_m}\to\mathcal A_m\to\nu_*\OO_P\to0$
has a surjective degree-zero residue map at $p_0$: the logarithmic
horizontal form $dq/q$ has residue one on every branch. Its kernel
is the one elementary modification
$f_{m*}\Omega^1=\OO_B(-2)$ of
$f_{m*}\mathcal A_m=\rho_*\OO(5)=\OO_B(-1)$.
The rationality of $P$ makes the higher direct images of the
residue quotient zero. Thus
$R^1f_{m*}\Omega^1=R^1f_{m*}\mathcal A_m$.
Since $H^1(A_m)=0$, its global sections have dimension
$4m+h^0(K_m)$, and Leray adds
$h^1(\OO_B(-2))=1$ in total degree one. This proves the formula
for $h^{1,1}$. Finally the sphere Chern numbers give
$\chi(\Omega^1)=-1$, and the vanishing just proved, together with
Serre duality, gives $h^{1,1}-h^{1,2}=1$. This proves (E.5).
\end{proof}

\begin{claim}{E2.3}{Open}{The remaining comb splitting and deformation calculations}
The integer $m$ and the stated limiting data do not determine in
this paper the splitting of the rank-three bundle $K_m$ in (E.6),
its exact number of sections, or the obstruction map of the comb
member. These parts of E.2 are unfinished. The four tangent
cohomology dimensions themselves are fully determined by (E.4).
\end{claim}
\begin{proof}[The remaining finite data and the computations made]
If $e_m$ is the length of $\operatorname{coker}\kappa_1$ in (E.6),
then its source has rank four and degree $-2$, its target is
$\OO_B(2m)$, and
\[
 \deg K_m=-2m-2+e_m.                                  \tag{E.7}
\]
The degree $-2$ follows from the two pairs of opposite stabilizer
characters of $E_m^*\otimes G_m$; its root determinant is trivial,
and each finite root point subtracts one from the coarse determinant.
Although (E.4a) determines the tangent dimensions independently of
the individual zeros of the Kodaira--Spencer determinant, the
support and elementary divisors of $Q_m$ still depend on that
period-section map. The comb Kuranishi problem further requires
its actual bracket and higher obstruction maps; the dimensions in
(E.4) do not determine those maps.

We evaluated the universal exact degree and dimension expressions
for $m=12,13,20,100$ in \texttt{d\_exact.py} and
\texttt{d\_comb.py}; the first case already
gives $h^{1,1}\ge49$, $h^{1,2}\ge48$, $h^0(T)=25$,
$h^1(T)=97$ and $h^2(T)=73$. The smallest undecided instance
is the splitting of (E.6) for one specified smoothed 12-tooth period
section, together with that member's obstruction map. The
attached construction supplies an existence argument and an open
set of smoothing choices, not the transition matrices of a uniquely
specified $X_{12}$. We have not chosen and solved its nonlinear
period-section smoothing equations, and have not run two numerical
methods on those missing matrices. Formulas (E.5)--(E.7) and the
growth bounds are proved, but they are not replacements for that
remaining calculation or for its Kuranishi bracket.
\end{proof}

\begin{claim}{E2.4}{Refuted}{The universal two-vector-field bound}
The assertion that every complex structure on $S^6$ has
$h^0(T)\le2$, printed in
\citet[Corollary 1.3; verified preprint p.~2]{HKP2000}
and repeated in \citet[Remarks 1.2 and 9.26]{AC2026}, is false
for the comb members.
\end{claim}
\begin{proof}
E2.2 gives $h^0(T_{X_m})=2m+1\ge25$ for $m\ge12$.
The source proof expresses a third field as a meromorphic linear
combination of two fields and invokes the Campana--Demailly--Peternell
claim that all meromorphic functions on a complex six-sphere are
constant. That implication is unavailable here: $f_m$ supplies
nonconstant meromorphic functions. The independently computed
$2m+1$ sections above give the counterexample, so this is a
refutation of the bound as well as a location of the failed step.
It does not refute a statement restricted to algebraic dimension
zero.
\end{proof}
